% 4_2_Simulation.tex
\subsection{Simulation Conditions}
We implemented a stochastic Susceptible-Infected-Recovered (SIR) model using a discrete-time chain binomial process.
The population consists of $N = 10,000$ individuals, sampled based on the Year 2000 European demographic distribution.
The transition rules are defined as follows:

\textbf{Infection ($S \to I$):}
At each daily time step, a susceptible individual $i$ becomes infected with a probability dependent on the number of their currently infected neighbors, $I_{\text{neighbors}, i}$.
The probability of infection is given by:
$$P(S_i \to I_i) = 1 - \exp(-\rho \cdot I_{\text{neighbors}, i})$$
where $\rho$ is the per-contact transmission probability.
To ensure comparability across networks with different mean degrees $\bar{k}$, $\rho$ is analytically derived from the basic reproduction number $R_0$ using the relationship $\rho = -\log(1 - R_0 / \bar{k}) / d$.
In our primary scenario, we set $R_0 = 2.5$.
This value is selected to reflect the transmission potential observed in the early stages of the COVID-19 pandemic and other significant respiratory outbreaks \parencite{alimohamadi2020estimate, billah2020reproductive}.

\textbf{Recovery ($I \to R$):}
Infected individuals recover at a constant daily rate $\gamma$, independent of the network structure.
The duration of infection is set to $d = 7$ days, yielding a recovery probability of:
$$P(I \to R) = \gamma = \frac{1}{d}.$$

Each simulation is initialized with 5 infected individuals ($I_0 = 5$).
To reflect realistic outbreak origins, the initial nodes in the structured networks (Model 1 and 2) are selected with probabilities weighted by their degree, while random selection is used for the homogeneous models.
This degree-weighted initialization accounts for the higher probability of an outbreak being introduced or first detected among highly sociable "hubs" in the population \parencite{lloyd-smith2005superspreading}.